\section{Generic groups}

So far we have treated the Fourier transform on the cyclic group $[\prod_{\catenumeratorin}\catdimof{\catenumerator}]$ using unfoldings.
More naturally, a tensor is defined as a complex field on the cartesian product group $\facstates$.
Let us now generalize to generic groups.

\subsection{Tensors as group fields}

A group field is a vector
\begin{align*}
    \hypercoreat{\catvariable} \wcols G \rightarrow \mathbb{C}
\end{align*}
where $\catvariable$ enumerated the $\cardof{G}=\catdim$ elements of the group.
The $\catindex$-th element of the group is then mapped to $\hypercoreat{\indexedcatvariable}$.

\begin{example}{Unfolding is not an isomorphism}
    Let us note that the unfolding of $[\prod_{\catenumeratorin}\catdimof{\catenumerator}]$ to $\facstates$ is not a group isomorphism.
    The group operation in $\prod_{\catenumeratorin}\catdimof{\catenumerator}$ is the integer summation
    \begin{align*}
        \catindex \circ \seccatindex =
        \catindex + \seccatindex \modspace \left(\prod_{\catenumeratorin}\catdimof{\catenumerator}\right) \, .
    \end{align*}
    In the unfolding the group operation is
    \begin{align*}
        (\catindexof{0},\ldots,\catindexof{\catorder-1}) \circ (\seccatindexof{0},\ldots,\seccatindexof{\catorder-1}=
        = (\catindexof{0}+\seccatindexof{0}\modspace \catdimof{0},\ldots,\catindexof{\catorder-1}+\seccatindexof{\catorder-1}\modspace \catdimof{\catorder-1}) \, .
    \end{align*}
    Intuitively, the difference of these groups is that in the unfolded group summation carry bits are set to $0$.

\end{example}

\subsection{Representations}

A representation is a map of $G$ to a matrix group.
Irreducible representations (irrep) $\rho$ are those, which are not equivalent to matrix groups in block form.
We represent them by tensors
\begin{align*}
    \rho[\catvariable,\cheadvariable{\insymbol},\cheadvariable{\outsymbol}] \, .
\end{align*}


\begin{example}{Cyclic group}\label{exa:cycGroupIrreps}
    For the cyclic group $G=[\catdim]$ the irreps are one-dimensional, that is $\headdim=1$, and enumerated by $\selindex\in[\catdim]$ as
    \begin{align*}
        \rho^{[\catdimof{\selindex}]}[\indexedcatvariable,\indexedselvariable,\cheadvariable{\insymbol}=0,\cheadvariable{\outsymbol}=0]
        = \expof{2\pi i \cdot \frac{\catindex \cdot \selindex}{\catdim}} \, .
    \end{align*}
    Note that for $\catdim=2$ this is the Hadamard matrix
    \begin{align*}
        \rho^{[2]}[\catvariable,\selvariable,\cheadvariable{\insymbol}=0,\cheadvariable{\outsymbol}=0]
        = \sqrt{2} \cdot \hgateat{\catvariable,\selvariable} \, .
    \end{align*}
\end{example}



\subsection{Finite Abelian groups}

Finite Abelian groups are equivalent to cartesian products of cyclic groups, that is
\begin{align*}
    G \sim \bigtimes_{\catenumeratorin}[\catdimof{\catenumerator}] \, .
\end{align*}
The irreps of finite Abelian groups are thus disconnected tensor networks
\begin{align*}
    \rho^{\bigtimes_{\catenumeratorin}[\catdimof{\catenumerator}]}[\shortcatvariables,\selvariableof{[\catorder]},\cheadvariableof{\insymbol}{[\catorder]},\cheadvariableof{\outsymbol}{[\catorder]}]
    = \bigotimes_{\catenumeratorin} \rho^{[\catdimof{\catenumerator}]}[\catvariableof{\catenumerator},\selvariableof{\catenumerator},\cheadvariableof{\insymbol}{\catenumerator},\cheadvariableof{\outsymbol}{\catenumerator}]
\end{align*}
where each $\rho^{[\catdimof{\catenumerator}]}[\catvariableof{\catenumerator},\selvariableof{\catenumerator},\cheadvariableof{\insymbol}{\catenumerator},\cheadvariableof{\outsymbol}{\catenumerator}]$ is the irrep tensor to $[\catdimof{\catenumerator}]$, see \exaref{exa:cycGroupIrreps}.


\subsection{Fourier transform}

The Fourier transform of a group field $\hypercoreat{\catvariable}$ maps each irrep $\rho$ to
\begin{align*}
    \sqrt{\frac{\headdim}{\cardof{G}}} \cdot \contractionof{\hypercoreat{\catvariable},\rho[\catvariable,\cheadvariable{\insymbol},\cheadvariable{\outsymbol}]}{\cheadvariable{\insymbol},\cheadvariable{\outsymbol}} \, .
\end{align*}

\begin{example}[Walsh-Hadamard transform]
    Consider the group $[2]^{\catorder}$, and a complex field on it by a quantum state $\qstateat{\shortcatvariables}$.
    The Fourier transform is then the application of Hadamard gates on each qubit, that is
    \begin{align*}
        \contractionof{
            \{\qstateat{\shortcatvariables}\}
            \cup \bigcup_{\catenumeratorin}
            \{\hgateat{\catvariableof{\catenumerator},\selvariableof{\catenumerator}}\}
        }{\selvariableof{[\catorder]}} \, .
    \end{align*}
    This is the Walsh-Hadamard transform of the quantum state.
\end{example}